% Part: computability
% Chapter: recursive-functions
% Section: pr-functions-computable

\documentclass[../../../include/open-logic-section]{subfiles}

\begin{document}

\olfileid{cmp}{rec}{cmp}
\olsection{ఆదిమ పునరావృత్త ప్రమేయాలు గణనీయమైనవి}

ప్రమేయం $h$ను కింది ఆదిమ పునరావృత్తి ద్వారా నిర్వచించామనుకుందాం:
\begin{eqnarray*}
h(\vec x, 0)     & = & f(\vec x) \\
h(\vec x, y + 1) & = & g(\vec x, y, h(\vec x, y))
\end{eqnarray*}
అలాగే $f$, $g$ ప్రమేయాలు గణనీయమైనవని అనుకుందాం. ($x_0$, \dots,
$x_{k-1}$లకు సంక్షిప్తంగా $\vec x$ను వాడుతున్నాం.) అప్పుడు
$h(\vec x,0)$ కేవలం $f(\vec x)$; అది గణనీయమని అనుకున్నాం కనుక
$h(\vec x,1)$ను స్పష్టంగానే గణించవచ్చు. $1=0+1$ కనుక
$h(\vec x,1)$ను కూడా గణించవచ్చు; అది
\begin{align*}
 h(\vec x, 1) & = g(\vec x, 0, h(\vec x, 0)) =  g(\vec x, 0, f(\vec x)).
\intertext{ఇదే విధంగా కొనసాగించి కింది వాటిని గణించవచ్చు:}
h(\vec x, 2) & = g(\vec x, 1, h(\vec x, 1)) = g(\vec x, 1, g(\vec x, 0, f(\vec x)))\\
h(\vec x, 3) & = g(\vec x, 2, h(\vec x, 2)) = g(\vec x, 2, g(\vec x, 1, g(\vec x, 0, f(\vec x))))\\
h(\vec x, 4) & = g(\vec x, 3, h(\vec x, 3)) = g(\vec x, 3, g(\vec x, 2, g(\vec x, 1, g(\vec x, 0, f(\vec x)))))\\
& \vdots
\end{align*}
కనుక సాధారణంగా $h(\vec x,y)$ను గణించడానికి, $h(\vec x,0)$,
$h(\vec x,1)$, \dots{}లను వరుసగా గణిస్తూ $h(\vec x,y)$ను చేరాలి.

ఈ విధంగా $f$, $g$ గణనీయమైతే, ఆదిమ పునరావృత్త నిర్వచనం ఒక కొత్త
గణనీయ ప్రమేయాన్ని ఇస్తుంది. $f$, $g_i$లు గణనీయమైతే, వాటి సంయుక్తం
కూడా గణనీయ ప్రమేయాన్నే ఇస్తుంది.

ప్రాథమిక ప్రమేయాలు $\Zero$, $\Succ$, $\Proj{n}{i}$ గణనీయమైనవి;
సంయుక్తం, ఆదిమ పునరావృత్తి గణనీయ ప్రమేయాల నుంచి గణనీయ ప్రమేయాలనే
ఇస్తాయి. అందువల్ల ప్రతి ఆదిమ పునరావృత్త ప్రమేయం గణనీయమైనది.

\end{document}
